\chapter{Implementación: arenas, RCU y motor de reglas}
\label{cap:implementacion}

Los caps.~\ref{cap:typegraph}--\ref{cap:dpo} fijaron el modelo formal:
qué tipos viven en el grafo, qué reglas DPO actúan sobre él. Este
capítulo baja a Rust: cómo el runtime aloja los nodos, cómo los lectores
ven snapshots consistentes mientras un escritor muta, y cómo cada
función POSIX en \citaCrate{sotfs-ops} desemboca en una mutación segura.

\section{Arenas: alocación dense con reuso}

\begin{defn}[Arena tipada]
\label{def:arena}
Una \emph{arena} \texttt{Arena<T, CAPACITY>} es un pool densamente
empaquetado de hasta \texttt{CAPACITY} valores de tipo \texttt{T}, con
un \emph{free-list} de slots reusables después de la liberación.
Operaciones: \texttt{alloc(v) -> Option<ArenaId>} (O(1)),
\texttt{free(id)} (O(1)), \texttt{get(id) -> Option<\&T>} (O(1)).
\end{defn}

El runtime usa una arena por tipo de nodo (inodes, dirs, caps, txns,
versions, blocks) y una arena para aristas. \texttt{CAPACITY} es
65\,536 para nodos y 131\,072 para aristas
(\citaLine{sotfs-graph/src/graph.rs}{31}). Los datos viven en
\texttt{Box<[MaybeUninit<T>]>} (heap-backed para evitar stack overflow:
una arena entera son ~35\,MiB).

\paragraph{Estructura interna.}

\begin{itemize}
  \item \texttt{data}: el array de slots \texttt{MaybeUninit<T>}.
  \item \texttt{state}: per-slot \texttt{Vacant} o \texttt{Occupied}.
  \item \texttt{free\_list}: stack de índices liberados, popea LIFO.
  \item \texttt{next\_bump}: índice del próximo slot bump-allocable.
\end{itemize}

\paragraph{Algoritmo de alocación.}

\begin{algorithm}[H]
\SetAlgoLined
\KwIn{Arena $A$, valor $v$ de tipo $T$}
\KwOut{\texttt{ArenaId} del slot ocupado, o \texttt{None} si llena}
\caption{\textsc{Arena.Alloc}}\label{alg:arena-alloc}

\eIf{$A.\textrm{free\_count} > 0$}{
  $A.\textrm{free\_count} \gets A.\textrm{free\_count} - 1$\;
  $\textit{slot} \gets A.\textrm{free\_list}[A.\textrm{free\_count}]$\;
}{
  \eIf{$A.\textrm{next\_bump} < \texttt{CAPACITY}$}{
    $\textit{slot} \gets A.\textrm{next\_bump}$\;
    $A.\textrm{next\_bump} \gets A.\textrm{next\_bump} + 1$\;
  }{
    \Return{\texttt{None}}\;
  }
}
$A.\textrm{data}[\textit{slot}] \gets v$\;
$A.\textrm{state}[\textit{slot}] \gets \texttt{Occupied}$\;
$A.\textrm{len} \gets A.\textrm{len} + 1$\;
\Return{\texttt{Some}(\texttt{ArenaId}($\sotxorigmathit{slot}$))}\;
\end{algorithm}

\begin{lema}[Densidad de la arena]
\label{lem:arena-densidad}
Tras una secuencia arbitraria de \textsc{Alloc}/\textsc{Free}, la
fracción de slots ocupados respecto del rango utilizado
$[0, \textrm{next\_bump})$ está acotada inferiormente por
$\textrm{len}/\textrm{next\_bump}$. La arena no fragmenta porque
\textsc{Free} siempre devuelve el slot al free-list y \textsc{Alloc}
prefiere ese pool antes de bump-alocar.
\end{lema}

\paragraph{Por qué importa.} \texttt{ArenaId} es \texttt{u32} (4 bytes)
y se mapea a un \texttt{u64} typed-id (\texttt{InodeId}, \texttt{DirId},
\dots) sumando 1 (los IDs en sotFS empiezan en 1; el slot 0 nunca se usa
como ID, sólo como sentinel). El ratio compacto evita cargar tablas
grandes en cache cuando se itera el grafo.

\section{Cómo se usa la arena en \texttt{TypeGraph}}

La \texttt{struct TypeGraph} (\citaLine{sotfs-graph/src/graph.rs}{52})
tiene siete arenas:

\begin{lstlisting}[language=rust]
pub struct TypeGraph {
    pub inodes: Arena<Inode, NODE_CAPACITY>,        // 65k
    pub dirs: Arena<Directory, NODE_CAPACITY>,
    pub caps: Arena<Capability, NODE_CAPACITY>,
    pub transactions: Arena<Transaction, NODE_CAPACITY>,
    pub versions: Arena<Version, NODE_CAPACITY>,
    pub blocks: Arena<Block, NODE_CAPACITY>,
    pub edges: Arena<Edge, EDGE_CAPACITY>,           // 131k
    // ... + índices secundarios BTreeMap
}
\end{lstlisting}

Una \texttt{TypeGraph} ocupa $\sim 35$\,MiB en heap por instancia
(siete arenas $\times 4$\,KiB de header + datos a demanda). Cap.~\ref{cap:transacciones}
discute por qué eso vuelve costoso \texttt{graph.clone()}.

\section{RCU: lectores lock-free, escritor único}

\begin{defn}[RcuGraph]
\label{def:rcu-graph}
\texttt{RcuGraph} mantiene \emph{dos} copias del \texttt{TypeGraph} en
heap (\texttt{Box<TypeGraph>}). Un índice atómico \texttt{active}
indica cuál ven los lectores; el escritor muta la inactiva, hace un
swap atómico, y espera a que los lectores que aún ven la copia vieja
terminen (\emph{rcu-synchronize}).
\end{defn}

El protocolo:

\begin{algorithm}[H]
\SetAlgoLined
\KwIn{$\textit{rcu}: \texttt{RcuGraph}$, mutador $f: \texttt{TypeGraph} \to ()$}
\caption{\textsc{RcuGraph.Write}}\label{alg:rcu-write}

\textsc{Acquire}($\textit{rcu}.\textrm{write\_lock}$)
$\textit{inactive} \gets 1 - \textit{rcu}.\textrm{active}.\textsc{Load}()$\;
$\textit{rcu}.\textrm{graphs}[\textit{inactive}] \gets \textit{rcu}.\textrm{graphs}[\textit{rcu}.\textrm{active}].\textsc{Clone}()$\;
$f(\textit{rcu}.\textrm{graphs}[\textit{inactive}])$\;
$\textit{rcu}.\textrm{active}.\textsc{Store}(\textit{inactive}, \texttt{Release})$\;
\textsc{RcuSynchronize}($\textit{rcu}$)
\textsc{Release}($\textit{rcu}.\textrm{write\_lock}$)\;
\end{algorithm}

\paragraph{Lectores.} Cada lector reclama un slot de
\texttt{reader\_epochs} (hasta 64 en paralelo, fijo en
\citaLine{sotfs-graph/src/rcu.rs}{39}), bumpea el global epoch y
escribe ese epoch en su slot. La sincronización del escritor consiste
en esperar a que todos los slots con epoch $\le$ swap-epoch retornen a
\texttt{INACTIVE}. \emph{No bloquea a nuevos lectores}: ellos ya verán
la copia nueva.

\begin{invariante}[RCU consistency]
\label{inv:rcu-consistency}
Un lector que entró antes del swap ve el grafo en su estado pre-swap
durante toda su lectura, aún si en paralelo otro escritor está mutando
la copia opuesta.
\end{invariante}

\paragraph{Trampa subtle.} \texttt{MAX\_READERS} fijo en 64; un
deployment con más lectores concurrentes que ese número entra a un
\emph{panic path} (\citaLine{sotfs-graph/src/rcu.rs}{31}). El número
viene del bump v0.2.1 cuando una FUSE daemon a 16 cores trivialmente
excedía los 8 originales. El fix definitivo (per-CPU counters o pool
dinámico) está en post-v0.3.

\section{El motor de reglas: cómo se aplica un paso DPO}

Recapitulando del cap.~\ref{cap:dpo}: el runtime no implementa un motor
DPO genérico. Cada operación POSIX es una función Rust que en orden
hace (i) matching, (ii) gluing condition, (iii) mutación. La estructura
canónica:

\begin{algorithm}[H]
\SetAlgoLined
\KwIn{$g: \&\texttt{mut TypeGraph}$, parámetros de la op}
\KwOut{\texttt{Result<X, GraphError>}}
\caption{\textsc{Esqueleto de cada op DPO en sotfs-ops}}\label{alg:op-skeleton}


$\textit{handles} \gets g.\textsc{Resolve}(\textit{params})$\;
\If{$\textit{handles} = \texttt{None}$}{\Return{\texttt{Err(NotFound)}}\;}


\If{$\neg\;\textsc{CheckPrecond}(g, \textit{handles}, \textit{params})$}{
  \Return{\texttt{Err(\textit{precond violada})}}\;
}


$\textit{newId} \gets g.\textsc{Alloc}(\dots)$\;
$g.\textsc{Insert}(\dots)$\;
$g.\textsc{UpdateIndices}(\dots)$\;
$g.\textsc{RecordProv}(\dots)$
\Return{\texttt{Ok}(\textit{newId})}\;
\end{algorithm}

\paragraph{Punto clave.} La mutación nunca toca un slot ocupado: o
aloca uno nuevo (vía \textsc{Arena.Alloc}) o libera uno viejo. Esto
preserva el \invref{inv:no-dangling} \emph{por construcción de la
arena}: un slot ya liberado no aparece en ningún índice secundario, y
los nuevos índices apuntan a slots recién alocados.

\section{Provenance log: el sidecar opcional}

Si la GTXN está corriendo bajo \fuse{} con
\texttt{SOTFS\_PROV\_SIDECAR} configurado
(cap.~\ref{cap:fuse}), cada op llama a
\texttt{record\_prov(op, inode\_id, detail)}
(\citaLine{sotfs-graph/src/graph.rs}{619}) que serializa una entrada
JSONL con timestamp, tipo de op, inode, cap-id (si presente) y dominio.
La política se inyecta vía \texttt{set\_cap\_ctx(CapContext)} antes de
cada op (\citaCommit{61e7faf}).

\section*{Resumen del capítulo}

\begin{itemize}
  \item Los nodos viven en arenas tipadas \texttt{Arena<T,
    CAPACITY>} con free-list LIFO; \textsc{Alloc}/\textsc{Free} son
    O(1) y la arena no fragmenta.
  \item \texttt{TypeGraph} aglutina siete arenas (seis de nodo + una de
    arista) más índices secundarios \texttt{BTreeMap}; ocupa
    $\sim 35$\,MiB en heap por instancia.
  \item \texttt{RcuGraph} mantiene dos copias y sincroniza a los
    lectores vía epoch counters; los escritores muta la inactiva y
    swappea con \texttt{Release}.
  \item Cada operación POSIX es una función Rust en \citaCrate{sotfs-ops}
    que sigue el esqueleto matching $\to$ gluing $\to$ mutación.
  \item Provenance log opcional, drainable, con cap-mediated context.
\end{itemize}

\section*{Ejercicios}

\begin{ejercicio}[\dificultad{$\star$}]
\label{ej:impl-1}
¿Cuál es \texttt{NODE\_CAPACITY} y \texttt{EDGE\_CAPACITY} en sotFS
v0.2.5?
\end{ejercicio}

\begin{ejercicio}[\dificultad{$\star$}]
\label{ej:impl-2}
¿Cuántos lectores concurrentes admite el RCU? ¿Qué pasa si se excede?
\end{ejercicio}

\begin{ejercicio}[\dificultad{$\star\star$}]
\label{ej:impl-3}
Calcule la ocupación de memoria de un \texttt{TypeGraph} vacío
(arenas vacías). Pista: \texttt{Box<[MaybeUninit<T>]>} aloca
\texttt{CAPACITY} entradas pero no las inicializa.
\end{ejercicio}

\begin{ejercicio}[\dificultad{$\star\star$}]
\label{ej:impl-4}
La función \texttt{rcu\_synchronize} es lineal en el número de
lectores. Diseñá un escenario donde un escritor podría \emph{starvarse}
si los lectores se renuevan más rápido que la sincronización. ¿Es
posible bajo carga real?
\end{ejercicio}

\begin{ejercicio}[\dificultad{$\star\star$}]
\label{ej:impl-5}
Mostrá que el algoritmo \textsc{Arena.Alloc} preserva la invariante
``\texttt{state[i] == Occupied} si y solo si \texttt{data[i]} está
inicializado''. Identifique los tres puntos de la prueba.
\end{ejercicio}

\begin{ejercicio}[\dificultad{$\star\star$}]
\label{ej:impl-6}
Diseñe un test de stress para la arena: $N$ alocaciones, $N/2$
liberaciones aleatorias, $N/2$ alocaciones adicionales. ¿Cuáles
invariantes verificar al final?
\end{ejercicio}

\begin{ejercicio}[\dificultad{$\star\star$}]
\label{ej:impl-7}
\trampref{tr:cycle-cost} habla de cómo un índice secundario
(\texttt{dir\_for\_inode\_idx}) elimina un O(N\textsuperscript{3}). En
términos de la arena: ¿podría implementarse el índice como otra arena
en lugar de \texttt{BTreeMap}? Discuta tradeoffs.
\end{ejercicio}

\begin{ejercicio}[\dificultad{$\star\star\star$}]
\label{ej:impl-8}
Reemplace los \texttt{reader\_epochs} fijos de 64 slots por un
mecanismo de pool dinámico per-CPU que escale a cualquier número de
lectores. Bosquejo el diseño en pseudo-Rust + invariantes que se
preservan.
\end{ejercicio}
